\documentclass[tensorial.tex]{subfiles}

\providecommand{\sL}{\mathcal{L}}
\providecommand{\sLp}{\mathcal{L}^\perp}

\begin{document}

\chapter{Estrutura Tensorial}\label{cap:tensorial}

\section*{Sumário do capítulo}

Este capítulo apresenta a álgebra hipercomplexa do ponto de vista
\emph{tensorial}. Em vez de partir da multiplicação $z \cdot w$ em $\R^n$,
começamos do objeto matemático mais primitivo: um tensor
$B \in V^{\otimes k}$ sobre um espaço vetorial $V$. A multiplicação
hipercomplexa emerge como o caso particular $k = 3$, com $V = \R^n$.
A tensorização $V^{\otimes k}$ para $k$ arbitrário é a generalização
natural; nela definimos o funcional de defeito
$\mathcal{E}_k$ cuja análise gradiente leva a constantes
combinatórias $a_k$, $b_k$ e $C_k^2$ explicitamente fechadas.

\section{Tensores e ações multilineares}\label{sec:tensor}

\begin{definicao}[Espaço tensorial]
Seja $V$ um espaço vetorial real de dimensão finita $\dim V = m$.
O \emph{$k$-ésimo produto tensorial} é
\[
  V^{\otimes k} \;=\; \underbrace{V \otimes V \otimes \cdots \otimes V}_{k \text{ fatores}}.
\]
Um \emph{$k$-tensor} é um elemento $B \in V^{\otimes k}$. Em coordenadas:
$B = \sum_{i_1, \ldots, i_k} B_{i_1 \cdots i_k}\, e_{i_1} \otimes \cdots \otimes e_{i_k}$.
\end{definicao}

\begin{definicao}[Ação multilinear]
Cada $B \in V^{\otimes k}$ define a aplicação multilinear
\begin{equation}\label{eq:acao-multilinear}
  \widetilde{B} : V^{k-1} \to V,
  \qquad \widetilde{B}(v_1, \ldots, v_{k-1}) := \sum_{i_1, \ldots, i_k} B_{i_1 \cdots i_k}\, v_1^{i_2} \cdots v_{k-1}^{i_k}\, e_{i_1}.
\end{equation}
\end{definicao}

\begin{observacao}
A correspondência $B \leftrightarrow \widetilde{B}$ é bijeção canônica
$V^{\otimes k} \cong \mathrm{Mult}(V^{k-1}, V)$. Distinguimos os dois
quando o contexto exigir; caso contrário escrevemos $B(\cdot)$ pela
ação direta.
\end{observacao}

\section{Multiplicação como tensor de ordem 3}\label{sec:mult-tensor}

\begin{definicao}[Multiplicação como tensor]
Uma \emph{multiplicação} em $V$ é um tensor $\mu \in V^{\otimes 3}$
(equivalente a uma forma bilinear $V \times V \to V$) via a ação
\eqref{eq:acao-multilinear} com $k = 3$:
\[
  z \cdot w := \widetilde{\mu}(z, w) = \sum_{i, j, l} \mu_{l i j}\, z^i w^j\, e_l.
\]
\end{definicao}

\begin{proposicao}[Decomposição simétrica/antisimétrica]
$\mu = \mu_S + \mu_A$, onde
\[
  \mu_S(z, w) := \tfrac{1}{2}(z \cdot w + w \cdot z), \qquad
  \mu_A(z, w) := \tfrac{1}{2}(z \cdot w - w \cdot z).
\]
$\mu_A$ está em $\Lambda^2 V \otimes V \subset V^{\otimes 3}$ (parte
antissimétrica nos primeiros dois fatores).
\end{proposicao}

\section{Decomposição canônica $V = \R \oplus V'$}\label{sec:decomp}

\begin{definicao}[Estrutura unitária]
Seja $V$ munido de produto interno $\langle \cdot, \cdot\rangle$ e um
elemento distinguido $\mathbf{1} \in V$ unitário. Define-se
\[
  V' := \mathbf{1}^\perp = \{v \in V : \langle v, \mathbf{1}\rangle = 0\},
  \qquad V \cong \R \oplus V'.
\]
Cada $z \in V$ se escreve $z = a_0 \mathbf{1} + \boldsymbol{a}$ com
$a_0 \in \R$ e $\boldsymbol{a} \in V'$.
\end{definicao}

\begin{teorema}[Forma canônica da multiplicação]\label{thm:forma-canon}
A classe de multiplicações $\mu \in V^{\otimes 3}$ que têm $\mathbf{1}$
como identidade ($\mathbf{1} \cdot z = z \cdot \mathbf{1} = z$) e
preservam decomposição
$z \cdot w = (\re) \mathbf{1} + \boldsymbol{c}$ via produto interno
$\langle \cdot, \cdot\rangle$ tem forma
\begin{equation}\label{eq:mult-canon}
  z \cdot w \;=\; \bigl(a_0 b_0 - \langle \boldsymbol{a}, \boldsymbol{b}\rangle\bigr) \mathbf{1}
              \;+\; \bigl(a_0 \boldsymbol{b} + b_0 \boldsymbol{a} + \boldsymbol{a} \times \boldsymbol{b}\bigr),
\end{equation}
onde $\times : V' \times V' \to V'$ é uma aplicação bilinear arbitrária.
\end{teorema}

\begin{proof}[Esboço]
A condição de identidade força que $\mu$, restrita aos blocos
$\R \otimes \R \to V$ e $\R \otimes V' \to V$, tenha exatamente os
coeficientes em \eqref{eq:mult-canon}. Os blocos $V' \otimes V' \to \R$ e
$V' \otimes V' \to V'$ são livres, dando o produto interno
$-\langle \boldsymbol{a}, \boldsymbol{b}\rangle$ (forçado pela compatibilidade
com $\mathbf{1}$ via conjugação) e $\boldsymbol{a} \times \boldsymbol{b}$.
\end{proof}

\begin{definicao}[Família $\mu_s$]
A sub-família com $\times_s := s \cdot \mathrm{cross}$ (onde $\mathrm{cross}$
é o produto vetorial canônico em $V' = \R^3$ ou $\R^7$) parametriza
álgebras hipercomplexas em $\R^4$ e $\R^8$ por $s \in \R$.
\end{definicao}

\section{Tensorização e ação em $V^{\otimes k}$}\label{sec:tensorizacao}

\begin{definicao}[Funcional de defeito]
Para $V = \R^m$ com produto interno standard e $B \in V^{\otimes k}$:
\begin{equation}\label{eq:E-k-tensor}
  \mathcal{E}_k(B) \;:=\; \mathbb{E}_{a_1, \ldots, a_{k-1} \sim \mathcal{N}(0, I_m)}\bigl[\bigl(\norm{B(a_1, \ldots, a_{k-1})}^2 - 1\bigr)^2\bigr].
\end{equation}
$\mathcal{E}_k$ é polinomial de grau $4$ em $B$, $O(m)$-invariante e
$S_{k-1}$-simétrico (pela permutação dos $a_i$'s).
\end{definicao}

\begin{observacao}[Recupera caso clássico]
Para $k = 3$ e $V = \R^n$ decomposto $\R \oplus \R^{n-1}$ via
\eqref{eq:mult-canon}, $\mathcal{E}_3$ mede o desvio da multiplicação
$\mu$ em ter norma multiplicativa: $\norm{z \cdot w}^2 = \norm{z}^2 \norm{w}^2$
(igualdade exata em álgebras de divisão normadas, Hurwitz).
\end{observacao}

\section{Subespaços $\sL_k$ e $\sLp_k$}\label{sec:L-perp}

\begin{definicao}[Subespaço $\sL_k$ (caso $k$ ímpar)]
Para $k \geq 3$ ímpar, $\sL_k \subset V^{\otimes k}$ é o subespaço gerado
pelas \emph{inclusões básicas}
\begin{equation}\label{eq:inclusao}
  \phi_{(s, P)}[i_1, \ldots, i_k] \;:=\; v_{i_s}\, \prod_{(a, b) \in P} \delta_{i_a, i_b},
\end{equation}
onde $s \in \{1, \ldots, k\}$ é o slot do vetor $v \in V$ e $P$ é
pareamento perfeito do conjunto $\{1, \ldots, k\} \setminus \{s\}$.
O número de tipos de inclusão é $N_k = k \cdot (k-2)!!$.

$\sLp_k$ é o complemento ortogonal em $V^{\otimes k}$ relativo ao
produto interno induzido.
\end{definicao}

\begin{proposicao}[Significado geométrico]
$\sL_k$ é o subespaço de tensores que ``contêm traços'': $B \in \sL_k$
se e somente se existe contração canônica de $B$ com produtos de $\delta$
de Kronecker que devolve um vetor não-nulo. $\sLp_k$ é o subespaço
\emph{traceless}: $\sum_i B_{\ldots, i, \ldots, i, \ldots} = 0$ para
toda contração de pares.
\end{proposicao}

\begin{definicao}[Caso $k$ par]
Para $k$ par, $\sL_k$ é gerado por inclusões com \emph{dois} vetores:
$\phi_{(s_1, s_2, P)}[i_1, \ldots, i_k] = v_{i_{s_1}} w_{i_{s_2}} \prod_{(a, b) \in P} \delta_{i_a, i_b}$,
com $N'_k = \binom{k}{2} (k-3)!!$ tipos.
\end{definicao}

\section{Constantes invariantes do programa}\label{sec:constantes}

A análise gradiente do funcional $\mathcal{E}_k$ produz três constantes
combinatórias com forma fechada explícita.

\begin{teorema}[Constante $a_k$]\label{thm:a-k-cap1}
Para $B \sim \mathcal{N}(0, I_{V^{\otimes k}})$ projetado em $\sLp_k$ e
$u_G$ a projeção de $\nabla \mathcal{E}_k(B)$ via uma inclusão básica fixa:
\begin{equation}
  \mathbb{E}[\norm{u_G}^2]_{\sLp} = a_k \cdot m^{(3k+1)/2} + O(m^{(3k-1)/2}),
  \qquad a_k = \begin{cases} 32 \cdot 6^{(k-1)/2} & k \text{ ímpar} \\ 96 \cdot 6^{(k-2)/2} & k \text{ par} \end{cases}.
\end{equation}
\end{teorema}

\begin{teorema}[Constante $b_k$]\label{thm:b-k-cap1}
Para $B \sim \mathcal{N}(0, I_{V^{\otimes k}})$ Gauss isotrópico:
\begin{equation}
  \mathbb{E}[\norm{\nabla \mathcal{E}_k(B)}^2] = 16 \cdot m^{3k} + O(m^{3k-2}),
  \qquad b_k = 16 \text{ universal}.
\end{equation}
\end{teorema}

\begin{teorema}[Constante leakage $C_k^2$]\label{thm:Ck-cap1}
A constante $C_k$ definida por
$C_k := \lim_{m \to \infty} m^{(2k-1)/2} \cdot \frac{\norm{P_{\sL_k}(\nabla \mathcal{E}_k(B))}}{\norm{\nabla \mathcal{E}_k(B)}}$
satisfaz, para $k$ ímpar:
\begin{equation}
  \boxed{C_k^2 \;=\; \frac{N_k \cdot a_k}{b_k} \;=\; 2 \cdot k \cdot (k-2)!! \cdot 6^{(k-1)/2}.}
\end{equation}
\end{teorema}

\begin{center}
\begin{tabular}{cccccc}
\toprule
$k$ & $N_k$ & $a_k$ & $b_k$ & $C_k^2$ & $C_k$ \\
\midrule
$3$ & $3$ & $192$ & $16$ & $36$ & $6$ \\
$5$ & $15$ & $1\,152$ & $16$ & $1\,080$ & $\sqrt{1080} \approx 32{,}86$ \\
$7$ & $105$ & $6\,912$ & $16$ & $45\,360$ & $\approx 213{,}0$ \\
$9$ & $945$ & $41\,472$ & $16$ & $2\,449\,440$ & $\approx 1565{,}1$ \\
\bottomrule
\end{tabular}
\end{center}

\section{Plano do livro}

A tensorialidade desta seção será o esqueleto sobre o qual:
\begin{itemize}[leftmargin=2em,nosep]
  \item Cap.~\ref{cap:matematica} (este): estrutura tensorial e constantes
        combinatórias.
  \item Cap.~\ref{cap:campos-locais}: caso $k = 3$, $V = \R^n$, decomposição
        em campos locais $\C_{\hat a}$, Hurwitz, mecânica algébrica.
  \item Cap.~\ref{cap:cohomologia}: cohomologia da família $\mu_s$,
        curvatura algébrica, conexão com $S^{n-2}$.
  \item Cap.~\ref{cap:fisica}: realização das estruturas tensoriais em
        eletromagnetismo, relatividade, mecânica quântica, fluidos.
  \item Cap.~\ref{cap:computacao}: aplicações em redes neurais, otimização
        combinatória, complexidade.
\end{itemize}

\section{Notações e convenções}

\begin{itemize}[leftmargin=2em,nosep]
  \item $V$, $W$: espaços vetoriais reais.
  \item $\dim V = m$: dimensão de $V$.
  \item $V^{\otimes k}$: $k$-ésimo produto tensorial.
  \item $B$, $G$: tensores em $V^{\otimes k}$ ou seus gradientes.
  \item $\R^n = \R \mathbf{1} \oplus V'$: decomposição canônica.
  \item $\mu, \mu_s$: multiplicação tensor / família parametrizada.
  \item $\times$, $\times_s = s \cdot \mathrm{cross}$: produto vetorial.
  \item $a_0 \in \R$, $\boldsymbol{a} \in V'$: parte real / imaginária.
  \item $\norm{\cdot}$: norma euclidiana.
  \item $\langle \cdot, \cdot\rangle$: produto interno.
  \item $\mathcal{E}_k$: funcional de defeito \eqref{eq:E-k-tensor}.
  \item $\sL_k$, $\sLp_k$: subespaços de tensores.
  \item $N_k$, $N'_k$: número de inclusões básicas (ímpar / par).
  \item $a_k$, $b_k$, $C_k^2$: constantes combinatórias.
\end{itemize}

\end{document}
